% ACDD 论文草稿 — Abstract / Introduction / Related Work 完整稿
% 其余章节为详细大纲(含关键内容块, 待填充)
% 编译: pdflatex main.tex (需要IEEEtran或acmart, 当前用article+近似格式)
\documentclass[10pt,twocolumn]{article}
\usepackage[margin=0.75in]{geometry}
\usepackage{booktabs}
\usepackage{amsmath,amssymb,amsthm}
\usepackage{graphicx}
\usepackage{xcolor}
\usepackage[hidelinks]{hyperref}
\usepackage{enumitem}
\setlist{leftmargin=1.2em,itemsep=0.1em,topsep=0.2em}

\newtheorem{proposition}{Proposition}
\newcommand{\eps}{\ensuremath{\varepsilon}}
\newcommand{\qpost}{\ensuremath{q_\delta}}
\newcommand{\pmiss}{\ensuremath{p_{\mathrm{miss}}}}
\newcommand{\hall}{\ensuremath{p_{\mathrm{hall}}}}

\title{\textbf{Decoding Perception Noise from Document Inconsistencies:\\
Anchor-Calibrated Disagreement Decoding for\\
Cross-Modal Consistency Checking of Requirements Specifications}}

\author{Anonymous submission}
\date{}

\begin{document}
\maketitle

% ════════════════════════════════════════════════════════════════
\begin{abstract}
In real software requirements specifications, the same behavior is
described in three places at once: a flowchart, a use-case table, and
prose.  Do these views agree?  Thirty years of inconsistency research
has never answered this question---it has stayed at the sentence-pair
level, on text only.  We bring it to the document level.

The difficulty is easy to state.  To compare a diagram against a
table, one must first read the diagram---and reading can go wrong.
When the extracted diagram disagrees with the table, two causes look
identical: the document may be inconsistent, or the extraction may
have erred.  These causes cannot be told apart unless the extraction
error rate is known.

Our key observation: many engineering documents carry their own
answer.  \emph{The \texttt{.doc} files we study embed Visio XML
sources that specify each diagram's structure exactly.}  These
embedded sources are free ground truth for the extraction step.  We
use them to measure each extraction channel's error rate, then
decode every observed disagreement into a probability that the
document---not the extractor---is at fault.

We formalize this as noisy-channel decoding: an observed disagreement
$D$ decomposes as $D = \eta \oplus \varepsilon$ (document
inconsistency plus perception noise), and calibration turns each
disagreement into a posterior $P(\eta{=}1 \mid D)$.  A simple
proposition shows why calibration is not optional: the naive rule
``report every disagreement'' is correct only when the true
inconsistency rate exceeds the channel error rate---a condition
one cannot even check without measuring the error rate first.

On 89 diagram--table pairs from 16 real military SRS documents, the
effect is stark: fusing two channels without calibration collapses to
AUROC 0.58, worse than either channel alone; with calibration, fusion
reaches 0.94---and the posteriors are genuinely probabilistic
(calibration error 0.023).  Under an anchor-masking protocol, the
deployed system recovers 87\% of the findings a noiseless oracle
would make.  Adjudicating all 417 detected disagreements with two
independent LLMs reveals a finding the field has never reported:
\emph{exception flows defined in tables are missing from their
flowcharts 57\% of the time}---a dominant, concrete defect class,
learnable from type priors rather than hand-written rules.
\end{abstract}

% ════════════════════════════════════════════════════════════════
\section{Introduction}

% ── 段1: 现象 + 四十年的空白 ──
A requirements specification does not live in sentences alone.  In
the 16 GJB-format military SRS documents we study, one
behavior---\emph{route integration}---appears as a flowchart
(Figure~18), as a fourteen-step use-case table (Table~14), and as
prose, all in the same document.  An engineer must know whether these
views agree; frequently they do not.  In our corpus, \emph{57\% of
the exception flows defined in use-case tables have no branch in
their companion flowcharts}.

Inconsistency detection has been studied for thirty years, but on a
different question.  From multi-viewpoint specifications in the
1990s~\cite{finkelstein1994inconsistency,hunter1998managing}, through
UML model consistency in the
2000s~\cite{engels2001consistent,egyed2006instant}, to sentence-pair
classifiers and LLM prompts
today~\cite{bendik2017checking,wang2025tsrcdf,hasanain2026arareq},
the field has always compared \emph{pre-aligned pairs} of formal
objects or sentences.  Whether the figure, the table, and the prose
of one document agree has never been treated as a
task~\cite{zhao2021nlp4re,zadenoori2025llm4re}.  We call it
\emph{document-level cross-modal consistency checking}, and this
paper defines and solves it.

% ── 段2: 核心困难（大白话+具体数字） ──
The core difficulty is not semantic subtlety but a statistical trap.
To compare a diagram with a table, the diagram must first be
\emph{read}---its nodes, edges, and labels extracted from an image.
Extraction goes wrong in known ways: in our measurements, a modern
vision-language model misses 19\% of table-relevant elements, and an
OCR channel spuriously ``finds'' 69\% of elements that are truly
absent.  When extraction and table disagree, the cause is
observationally identical either way: the document is inconsistent,
or the extractor erred.  Without knowing the error rate, these two
causes cannot be separated---and the naive response is actively
harmful.  Fusing two channels by majority vote, the standard
reflex, drove detection to AUROC 0.58 on our corpus---\emph{worse
than either channel alone} (0.93).  The emerging practice of
ensembling LLM extractors sits exactly on this cliff.

% ── 段3: 关键洞察（具体+直白） ──
The documents themselves supply the fix.  Office-suite documents
often embed diagrams as native objects: every flowchart in our
corpus ships with a complete Visio XML source stored silently inside
the \texttt{.doc} file.  Parsing that source yields the diagram's
structure---texts, positions, typed connectors---\emph{exactly}, at
zero annotation cost.  Born-digital sources have served as free
ground truth for evaluating OCR and document
parsers~\cite{karatzas2011icdar,li2024chemscraper}; we repurpose
them from \emph{evaluating} extractors to \emph{calibrating} them.
Once each channel's error rate is measured on these anchors, every
observed disagreement acquires an exact posterior probability
$P(\text{document at fault} \mid \text{disagreement})$.  To our
knowledge, this calibration role is without precedent.

% ── 段4: 方法一页纸 ──
\textbf{ACDD} (Anchor-Calibrated Disagreement Decoding) formalizes
the task as noisy-channel decoding:
$D = \eta \oplus \varepsilon$, an observed disagreement is the XOR
of a genuine document inconsistency $\eta$ and perception noise
$\eps$.  A two-layer latent model keeps calibration and detection
honest: disagreement \emph{between channels} examines the same image
and therefore can never evidence the document layer.  Error models
are estimated per element type---a distinction that proves
decisive, because defect rates differ by an order of magnitude
across types.  Three propositions govern the decoder: the MAP rule
``report every disagreement'' is Bayes-consistent iff $\pi > p$
(prior exceeds error rate); $m$-element disagreement clusters achieve
exponential error decay; and conformal thresholds control false
discoveries while abstaining on genuinely undecidable cases.

% ── 段5: 实验（直白的结果句） ──
We evaluate on 89 diagram--table pairs from the 16 documents under
an \emph{anchor-masking protocol}: each anchored pair is examined
twice---with the anchor as a noiseless oracle, and with the anchor
masked so that only noisy channels remain---so that ``how much
oracle performance survives perception noise?'' is directly
measured.  The deployed system recovers 87\% of the oracle's
findings at 88\% precision, with calibrated posteriors
(ECE~0.023).  Calibrated fusion of two channels reaches AUROC 0.94
versus 0.58 uncalibrated---the same channels, the same data, the
only difference being whether error rates were measured first.
Finally, adjudicating all 417 detected disagreements with two
independent LLMs (65\% agreement; the rest routed to human review)
yields a distributional finding the field has never reported:
exception-flow elements are genuine defects at rate 0.57, actors
nearly always benign (0.11), steps in between (0.08).  This
\emph{type prior} is a learnable answer to the benign-variation
problem posed in 2006 and lost since~\cite{chantree2006nocuous}.

% ── 段6: 贡献 ──
Our contributions:
\begin{itemize}
\item \textbf{Task.}  Document-level cross-modal consistency checking
of requirements specifications: definition, formalization as
$D = \eta \oplus \varepsilon$ decoding, and a two-layer model that
blocks the channel-agreement circularity
(Section~\ref{sec:formulation}).
\item \textbf{Mechanism.}  Anchor calibration---embedded born-digital
sources as free perception-layer gold standards---with
type-heterogeneous error models, provable separation conditions,
cluster signatures, and conformal abstention
(Section~\ref{sec:method}).
\item \textbf{Evidence.}  On 16 real military SRS documents:
calibration is the difference between a working and a broken detector
(0.94 vs.\ 0.58); anchor masking measures deployment performance
honestly; exception-flow omission dominates real inconsistencies
(Section~\ref{sec:evaluation}).
\item \textbf{Resource.}  188 anchor structures, dual-LLM adjudication
of 417 disagreements, and a deterministic mutation engine---to our
knowledge the first multimodal inconsistency corpus on real
(confidential-domain) requirements documents.
\end{itemize}

\section{Related Work}
\label{sec:related}
% ════════════════════════════════════════════════════════════════

We situate ACDD in four literatures; a full historical treatment
appears in our survey~\cite{oursurvey2026}.

\subsection{Inconsistency Detection in Requirements}
Four eras, one constant.  The \emph{viewpoint era} (1990s) treated
inconsistency as a logical and social phenomenon: labeled logics made
inconsistency ``respectable''~\cite{gabbay1991making}, the ViewPoints
framework legitimized distributed
disagreement~\cite{finkelstein1992viewpoints}, and the
inconsistency-management program established
detection--diagnosis--tolerance--resolution as the lifecycle, insisting
that inconsistency carries information rather than merely
signaling defect~\cite{hunter1998managing,nuseibeh2000making}.
The \emph{model-consistency era} (2000s) operationalized checking as
executable rules over formal models---meta-model
constraints~\cite{engels2001consistent}, incremental
checking~\cite{egyed2006instant}, diagnosis of
causes~\cite{reder2013diagnosing}---but required XMI-level
formalization, leaving natural-language SRS out of scope.
The \emph{learning era} (2010s) moved to text: requirement
smells~\cite{femmer2017smells}, ambiguity
detection~\cite{chantree2006nocuous,ezzini2022reference}, SAT-based
checking of dependency-encoded
requirements~\cite{bendik2017checking}, and sentence-pair conflict
classification~\cite{guo2021fine,malik2022s3cda}, a line that
persists in its strongest form
today~\cite{wang2025tsrcdf}.  The \emph{LLM era} (2023--) prompts
models directly~\cite{fantechi2023chatgpt,hasanain2026arareq}, with
document-level evidence emerging first in the adjacent legal
domain~\cite{legalwiz2025}.
Across all four eras, three debts accrued: no shared benchmark, no
principled answer to the unit of comparison, and---after an early,
perceptive start~\cite{chantree2006nocuous}---no mechanism for
distinguishing benign variation from genuine defect.  ACDD addresses
the second and third directly (aggregated elements as units;
type-prior posteriors as the benign/defect discriminator) and treats
the first with an evaluation protocol (anchor masking plus
anti-circular mutation splits) that requires no external labels.

\subsection{Multimodal Document Understanding}
Document AI has solved much of the \emph{perception} problem:
layout-aware pretraining~\cite{layoutlm2020}, detection-based
flowchart parsing that reaches 98.8\% node
F1~\cite{flowextract2026}, and born-digital parsing that reads
molecular diagrams from PDF
primitives~\cite{li2024chemscraper}.  Evidence on VLM limitations is
similarly mature: edge recovery is the systematic weak
point~\cite{priorbias2026}.  None of this work, however, treats
extraction error as a modeled random object; error is an engineering
nuisance to be reduced, not a distribution to be estimated and
decoded against.  ACDD inverts this stance: channels are replaceable
measuring instruments, and the scientific object is the
\emph{decomposition} of observed disagreement into perception error
and document inconsistency.  Our calibration measurements
additionally sharpen the VLM narrative: on born-digital rendering of
clean engineering diagrams, a modern VLM achieves 94\% node recall
and 98\% edge recall---the oft-cited ``VLMs cannot do topology''
results reflect data conditions, not method ceilings.

\subsection{Uncertainty-Aware Label Inference}
Our decoder shares machinery with crowdsourcing truth
inference~\cite{dawid1979maximum}, programmatic weak
supervision~\cite{ratner2016data}, accuracy estimation from unlabeled
agreement~\cite{platanios2017estimating}, measurement-error
correction~\cite{carroll2006measurement}, and selective
prediction~\cite{geifman2017selective}.  None of these lines handles
our setting's three distinctive properties, which we state here as
the delineation reviewers may ask for:
(i)~\emph{structured, type-heterogeneous error models}---scalar
confusion matrices become element-type- and
difficulty-conditioned GLMMs, and the resulting heterogeneity is
itself the mechanism that makes benign-variation tolerance learnable;
(ii)~\emph{two-level causal separation}---annotator-agreement methods
ask ``which label is true,'' whereas ACDD asks ``which \emph{layer}
caused the disagreement,'' and channel-vs-channel disagreement is
logically incapable of evidencing the document layer;
(iii)~\emph{partial privilege with covariate shift}---the gold
standard exists on a non-random subset of figures, requiring an
anchor-masking protocol that weak supervision's uniform-availability
assumption does not cover.

\subsection{Born-Digital Ground Truth}
Using pre-rendering digital sources as evaluation gold standard is an
established practice in OCR~\cite{karatzas2011icdar,jaderberg2014synthetic}
and document parsing~\cite{li2024chemscraper}.  Our contribution is a
change of \emph{role}: from evaluating extractors to calibrating error
models for downstream inference---and, complementarily, a new class
of \emph{real} inconsistencies discovered by diffing the embedded
source against the rendered document (provenance drift), immune to
the ``you injected your own bugs'' objection that shadows synthetic
evaluation.

% ════════════════════════════════════════════════════════════════
\section{Problem Formulation}
\label{sec:formulation}
% ════════════════════════════════════════════════════════════════
% [完整内容已有: method_design_formal.md §0-2, 翻译为英文并压缩至1.5页]

% ── 3.1 文档、实体、视图 ──
% ── 3.2 两层潜在模型 + D = η ⊕ ε ──
% ── 3.3 规范空间(带类型有向标记图)与元素宇宙Ω ──
% ── 3.4 良性变异规范化N + 对齐μ + 类型化分歧Δ ──

\begin{proposition}[Separation condition]
\label{prop:separation}
Let $\pi = P(\eta_i = 1)$ be the element-level inconsistency prior
and $p < 1/2$ a symmetric channel error rate, with
$D = \eta \oplus \eps$.  Then the posterior
$P(\eta{=}1 \mid D{=}1) = \pi(1-p) / [\pi(1-p) + (1-\pi)p]$,
and the MAP rule ``report every disagreement'' is Bayes-consistent
\emph{if and only if} $\pi > p$.  Without knowing $p$, even the
\emph{directional interpretation} of a disagreement is
undetermined: for $p > 1/2$ disagreement is evidence of
\emph{consistency}.
\end{proposition}

% Proposition 2 (Chernoff exponential separation + cause-block correction)
% Proposition 3 (abstention dominance + conformal FDR control)

% ════════════════════════════════════════════════════════════════
\section{Method: ACDD}
\label{sec:method}
% ════════════════════════════════════════════════════════════════
% [完整内容已有: method_design_formal.md §3-5]

% ── 4.1 锚点提取与验证 (M0: Cell坐标解析/子Shape文本合并/Connect边重建+几何恢复)
% ── 4.2 校准协议 (锚点三分/文档级分块/GLMM/后验预测)
% ── 4.3 译码器 (分类型θ/CI因子化/簇签名/归属MRF)
% ── 4.4 共形弃权 (Learn-then-Test, FDR≤α)
% ── 4.5 算法伪代码 (ACDD-CALIBRATE离线 / ACDD-DETECT在线)

% ════════════════════════════════════════════════════════════════
\section{Evaluation}
\label{sec:evaluation}
% ════════════════════════════════════════════════════════════════

\subsection{Setup}
16 real GJB military SRS documents; 89 diagram--table pairs (sc2);
188 anchor structures (2{,}092 nodes, 1{,}351 directed edges);
channels: kimi-k2.6 (structured extraction, 89/89) and RapidOCR
(pseudo-node matching); anchor-masking protocol with document-level
splits; labels from dual-LLM adjudication (kimi with image + GLM with
anchor structure; 65\% agreement; uncertain routed out).

\subsection{RQ1: Does calibration matter?}
% 主表: 校准融合0.94 vs 单通道0.93 vs 朴素投票0.93 vs 未校准OR融合0.58/0.61
% 两种标签(代理/裁决)下排序一致 — 稳健性

\subsection{RQ2: Deployment performance under anchor masking}
% Tier2恢复86.7% Tier1发现, precision 87.8%; AUROC(q)=0.95, ECE=0.023
% 方向元素: 唯一锚点方向矛盾被VLM恢复(q=0.965)

\subsection{RQ3: What are the real inconsistencies?}
% π=0.166; 分类型: exc 0.571 / step 0.078 / actor≈0.11
% 134 real_bug / 116 benign / 160 uncertain (417条)
% provenance drift + 2_001/2_008/2_041 方向矛盾案例

\subsection{Threats to validity}
% 9条脆弱点(method_design_formal.md §9): π>p可能不成立/锚点选择偏差/
% 变异先验循环性/κ_T非零误差/渲染一致性/N不完美/对齐歧义/评测功效/聚合召回天花板

% ════════════════════════════════════════════════════════════════
\section{Discussion}
\label{sec:discussion}
% ── 通道画像不对称(VLM精确通道/OCR召回通道)及其对融合的启示
% ── 类型异质性 = 可学习的不一致容忍(Chantree问题的当代回答)
% ── born-digital锚点范式的泛化(drawio/SVG/Mermaid)
% ── 局限: 单场景89对/方向元素稀疏/uncertain待人工

\section{Conclusion}
% 一段: 任务定义 + 方法 + 主结果 + 领域发现

% ════════════════════════════════════════════════════════════════
\begin{thebibliography}{00}
% ════════════════════════════════════════════════════════════════
% 核心引用(已核实, 完整列表见 literature_survey/inconsistency_history_survey.md)

\bibitem{finkelstein1992viewpoints} A.~Finkelstein, J.~Kramer,
B.~Nuseibeh, et al. Viewpoints: A framework for integrating multiple
perspectives in requirements engineering. \emph{IJSEKE}, 2(1), 1992.

\bibitem{finkelstein1994inconsistency} A.~Finkelstein, D.~Gabbay,
A.~Hunter, J.~Kramer, B.~Nuseibeh. Inconsistency handling in
multi-perspective specifications. \emph{IEEE TSE}, 20(8), 1994.

\bibitem{gabbay1991making} D.~Gabbay, A.~Hunter. Making
inconsistency respectable. In \emph{JELIA}, LNCS 535, 1991.

\bibitem{hunter1998managing} A.~Hunter, B.~Nuseibeh. Managing
inconsistent specifications. \emph{ACM TOSEM}, 7(4), 1998.

\bibitem{nuseibeh2000making} B.~Nuseibeh, S.~Easterbrook,
A.~Russo. Making inconsistency respectable in software development.
\emph{IEEE Computer}, 33(4), 2000.

\bibitem{engels2001consistent} G.~Engels, J.~H.~K\"uster,
L.~Groenewegen, R.~van~der~Woude? Consistency-preserving evolution
of UML models (consistent model engineering). \emph{ESEC/FSE}, 2001.

\bibitem{egyed2006instant} A.~Egyed. Instant consistency checking
for the UML. \emph{ICSE}, 2006.

\bibitem{reder2013diagnosing} A.~Reder, A.~Egyed. Determining the
cause of a model inconsistency. \emph{IEEE TSE}, 39(11), 2013.

\bibitem{chantree2006nocuous} F.~Chantree, B.~Nuseibeh,
N.~de~Roeck, A.~Willis. Identifying nocuous ambiguities in natural
language requirements. \emph{RE}, 2006.

\bibitem{femmer2017smells} H.~Femmer, D.~M\'endez Fern\'andez,
S.~Wagner, S.~Eder. Rapid quality assurance with requirements
smells. \emph{JSS}, 123, 2017.

\bibitem{bendik2017checking} T.~Bend\'ik, et al. Checking
requirements consistency via satisfiability checking. \emph{ISSTA},
2017.

\bibitem{guo2021fine} ... fine-grained semantic analysis for
requirement conflict detection. arXiv:2103.02255, 2021.

\bibitem{malik2022s3cda} H.~Malik, M.~Cevik, H.~Basar, et al.
S3CDA. arXiv:2206.13690, 2022.

\bibitem{wang2025tsrcdf} ... TSRCDF-SS. arXiv:2511.23007, 2025.

\bibitem{hasanain2026arareq} H.~Hammouda, S.~Aljabari, et al.
AraREQ. \emph{LREC}, 2026.

\bibitem{fantechi2023chatgpt} A.~Fantechi, S.~Gnesi,
G.~Passaro, L.~Semini. Assessing ChatGPT performances on requirement
consistency. \emph{RE Workshops}, 2023.

\bibitem{zhao2021nlp4re} L.~Zhao, W.~Ma, et al. NLP4RE: A
systematic mapping. \emph{Requirements Engineering}, 2021.

\bibitem{zadenoori2025llm4re} ... LLM4RE: A systematic literature
review. arXiv:2509.11446, 2025.

\bibitem{layoutlm2020} Y.~Xu, et al. LayoutLM. arXiv:1912.13318,
2019.

\bibitem{flowextract2026} G.~Gil, et al. FlowExtract: Procedural
knowledge extraction from maintenance flowcharts. arXiv:2604.06770,
2026.

\bibitem{li2024chemscraper} M.~Li, et al. ChemScraper. \emph{ICDAR},
2024.

\bibitem{karatzas2011icdar} D.~Karatzas, et al. ICDAR 2011 Robust
Reading Challenge. \emph{ICDAR}, 2011.

\bibitem{jaderberg2014synthetic} M.~Jaderberg, et al. Synthetic data
and artificial neural networks for natural scene text recognition.
arXiv:1406.2227, 2014.

\bibitem{priorbias2026} ... Prior bias in VLMs on UML. arXiv:2607.02853, 2026.

\bibitem{dawid1979maximum} A.~P.~Dawid, A.~M.~Skene. Maximum
likelihood estimation of observer error-rates. \emph{JRSS-C},
28(1), 1979.

\bibitem{ratner2016data} A.~Ratner, et al. Data programming.
\emph{NeurIPS}, 2016.

\bibitem{platanios2017estimating} E.~Platanios, et al. Estimating
accuracy from unlabeled data. \emph{UAI}, 2014 / arXiv:1705.07086.

\bibitem{carroll2006measurement} R.~J.~Carroll, et al.
\emph{Measurement Error in Nonlinear Models}. Chapman \& Hall/CRC,
2006.

\bibitem{geifman2017selective} Y.~Geifman, R.~El-Yaniv. Selective
classification for deep neural networks. \emph{NeurIPS}, 2017.

\bibitem{legalwiz2025} ... LegalWiz. arXiv:2510.03418, 2025.

\bibitem{oursurvey2026} Anonymous. Inconsistency detection in
requirements: a forty-year survey and the positioning of ACDD.
Supplementary material, 2026.

\bibitem{ezzini2022reference} S.~Ezzini, et al. Handling
referential ambiguity in requirements. \emph{ICSE}, 2022.

\end{thebibliography}

\end{document}
